\section{Conclusions and Future Work}
% We have uncovered a mechanism which leverages a straightforward sampling and voting that is unsupervised, simplistic, and universally applicable, capable of yielding performance akin to that of sophisticated frameworks. This mechanism can operate orthogonally to other prompt engineering techniques. We have delved into the principles and the scope of applicability of the mechanism, and have extended it in an unsupervised manner to achieve superior results. Furthermore, it supports the application of heterogeneous models in a more cost-effective fashion. 
% In conclusion, our pioneering research presents the first comprehensive study of the "More Agents" mechanism, substantiating its capability to significantly boost the performance of various LLMs in diverse settings enabling even less advanced LLMs to match or surpass their sophisticated counterparts. Exploring the mechanism's intrinsic properties has yielded insights that inform its applicability and have facilitated the refinement of its effectiveness. Our analyses yield a nuanced comprehension of the mechanism's functionality and optimization potential for various LLM applications.  Our findings offer a valuable framework for enhancing LLM strategies and signal exciting avenues for future exploration.

% In this paper, we conduct a comprehensive investigation into a phenomenon: the enhancement of LLM performance with an increase in the number of agents or samplings. 
In this paper, we report that \textit{more agents is all you need}, i.e., simply adding more instantiated LLM agents is what you need to obtain a better LLM performance in processing complex tasks, without bothering complicated methods, such as CoT pipelines, multi-agent collaboration frameworks, etc. 
We have conducted the first comprehensive study in the literature to understand such a ``scaling law'', including when it holds and how to facilitate its occurrence. 
%Specifically, we propose a simple sampling-and-voting method and conduct a comprehensive study to verify its generalizability. Experiments are conducted on a variety of datasets, including arithmetic reasoning, general reasoning, and code generation tasks, using various LLMs. 

%The results indicate that our simple sampling-and-voting method for instantiating agents can generally improve the performance of LLMs by increasing the ensemble size. 
The results indicate that Agent Forest can generally improve the performance of LLMs by increasing the ensemble size. 
Importantly, this method is orthogonal to different existing methods, which can lead to further improvements when combined with them. 

Furthermore, we observe that the performance gains are influenced by the difficulty of the task. To explore this correlation, 
we isolate and analyze three dimensions of task difficulty: the inherent difficulty, the length of reasoning steps, and the prior probability of a correct answer. 
We find that: 1) the performance gains increase then decrease by rising the inherent difficulty; 2) performance gains increase with the number of steps;  and 3) performance increases with the prior probability. Based on these properties, we also develop ways to boost the effectiveness of ``More Agents''. 

Considering that each input remains the same when we increase the number of agents, the sampling phase can be optimized to reduce the cost. Nonetheless, such a challenge of escalating costs commonly exists in works requiring multiple LLM calls \cite{cotsc, debate}. This aligns with recent findings as discussed in \cite{betterbenchmark} that current benchmarks focus narrowly on accuracy, leading to costly agents. Additionally, the lack of adequate holdout sets and standardization in evaluation practices further complicates the development of cost-effective agents. Addressing these issues can enhance the practical efficiency of AI agents in real-world scenarios. %For example, the CoT-SC method, which enhances reasoning performance through the generation of multiple CoTs, is similarly impacted by cost considerations. Likewise, in the Debate framework, cost issues are exacerbated as the agent number grows, due to the communication mechanism designed. 
We leave it as a future work to optimize. % However, this is actually a common issue presented not only in our method but also in other methods. For instance, CoT-SC improves reasoning performance by generating multiple though chains, which is bothered by cost issue. Similarly, within the Debate framework, as the number of agents increases, the cost is amplified due to the debate communication mechanism
%Our analysis is based on specific arithmetic tasks, which are considered a simplified form of general tasks. We plan to conduct analyses on more general tasks in future work. Also, 

% Our analysis based on the premise of structured queries, specifically those expanding upon arithmetic tasks. For generative queries, as well as open-domain question-answering scenarios that lack an explicit structured question format, discussions have not been elaborated.
%Our findings contribute to the growing body of knowledge on LLM optimization and pave the way for more sophisticated applications of these models in complex problem-solving scenarios.


%In conclusion, we've developed a sampling and voting that works independently or in tandem with other methods. Our comprehensive experiments across LLMs and datasets confirm its broad applicability.  Our analysis revealed that gains are tied to task difficulty, leading us to identify three key properties that we used to further improve our method.

% \section{Future Works}
% Our analysis unfolds on the premise of structured queries, specifically those expanding upon arithmetic problems. For generative inquiries, as well as open-domain question-answering scenarios that lack an explicit structured question format, discussions have not been elaborated. Our aspiration lies in furnishing a lucid analysis through structured methodologies and imparting insights for forthcoming related endeavors.

% \section{Social Impact}
% This paper presents work whose goal is to advance the field of Machine Learning.  There are many potential societal consequences of our work, none which we feel must be specifically highlighted here.